\chapter{结语：从空间感知到空间智能}
\label{ch:epilogue}

% ════════════════════════════════════════════════════════════════
%  新部「学习时代的 SLAM 与空间 AI」收官章，也是全书「感知与空间智能」主线
%  之末、转向「规划与控制」之前的枢纽。
%  反思性、相对简短（远短于技术章）：回望全书 SLAM 弧线（经典几何 → 概率估计
%  → 学习赋能 → 开放世界空间 AI），提炼贯穿全书的少数核心思想，点出尚未解决
%  的根本问题，把读者自然过渡到后续规划与控制。
%  吸收 SLAM Handbook ch19《Epilogue》——它是若干资深学者的格言集。
%  最高准则（也是全书最高 ventriloquize 风险）：绝不照搬、绝不逐条「某人说」，
%  把其精神熔成本书自己的收束口吻，事实与多方观点用本书语气综述 + \cite。
%  缝合：首尾呼应绪论（ch:intro / sec:intro-history / sec:intro-need /
%  sec:intro-frontier）与前言对「经典→学习→具身」的承诺；回望本部四章
%  （ch:learning / ch:dynamic / ch:semantic / ch:openworld）与更早的经典 SLAM；
%  提炼「同一张因子图、同一条 MAP 母方程 eq:intro-nls 贯穿始终」的主线；
%  向后过渡到 ch:planning / ch:control。
%  本章精简：不设前置自测/学习目标/知识点总表/习题，改用「开放问题/展望/思考题」。
% ════════════════════════════════════════════════════════════════

\begin{quote}
\small\itshape
我们从一句最朴素的追问出发——一台机器，置身于一个它一无所知的世界，能否一边走、一边把这世界画下来，同时又知道自己身在画中的何处？走到这里，是时候合上地图，回望来路了。
\end{quote}

\noindent 全书行至此处，SLAM 这条主线已经走完。从 \cref{ch:intro} 第一次把“同时定位与建图”写成一个“鸡生蛋”的循环（\cref{sec:intro-chicken-egg}），到刚刚结束的这一部里网络一次前向就吐出三维点图——三十余年的演进，被我们压进了不到一半篇幅。这一章不再有新的公式要推、新的系统要拆；它要做的，是\emph{回望}：我们究竟走了一条怎样的弧线，沿途有哪几个思想一以贯之地撑着这条弧线，哪些根本问题至今仍悬而未决，以及——为什么仅有“感知与建图”还远远不够，机器人还必须学会\emph{决策与行动}，从而把读者引向全书的后半程。

% ════════════════════════════════════════════════════════════════
\section{回望一条弧线：从几何到智能}
\label{sec:epilogue-arc}

\paragraph{我们走过的四步。}
若把全书 SLAM 主线抽成一条最简的曲线，它大致经过四个驿站，每一步都没有推翻前一步，而是把前一步\emph{包进}更大的图景里：

\begin{enumerate}
  \item \textbf{经典几何}。最初的问题是纯几何的：相机怎样成像、两视之间的对极约束、如何由对应三角化出三维点、如何用光束法平差把这些点与位姿一致地拟合（\cref{ch:camera,ch:vo,ch:nlopt}）。这一步教会我们\emph{世界是有结构的}，而结构可以被刚体运动与投影几何精确刻画。
  \item \textbf{概率估计}。几何给出“应该看到什么”，但传感器永远带噪声、世界永远不完全可知。于是 SLAM 被重述为一个\emph{状态估计}问题（\cref{def:intro-state-est}）：不是求一个确定的答案，而是求一个\emph{带不确定性}的后验。滤波与优化两条路（\cref{sec:intro-two-views}）由此分野又合流，最终汇成贯穿全书后端的那条 MAP 母方程（\cref{eq:intro-nls}）。这一步教会我们\emph{要对自己的无知诚实}，并把这份无知量化成协方差。
  \item \textbf{学习赋能}。手工写下的残差、权重与对应关系，终究受限于人对世界的先验想象。当数据足够多、网络足够强，我们把母方程里这些“手工件”逐个交给网络去学（\cref{ch:learning}）：网络出更鲁棒的特征与光流、出度量深度以钉住单目尺度、乃至把整个优化层做成可反向传播的一层。这一步教会我们\emph{几何不必被抛弃，而可以被学习\emph{包裹}}——让网络补它擅长的（从像素里榨证据），让几何做它擅长的（用刚体约束把证据一致地融成轨迹与地图）。
  \item \textbf{开放世界的空间 AI}。最后，地图本身的“语言”被拓宽了：从只记几何，到记“这是什么”（\cref{ch:semantic} 的度量-语义），到在动态与可变形的真实世界里维持一致（\cref{ch:dynamic}），再到用基础模型的开放词汇嵌入去回答\emph{任何}概念的查询（\cref{ch:openworld}）。这一步教会我们\emph{地图不只是给机器人“看路”的，更是给它“理解”与“被询问”的}——它开始承载语义、面向任务、连接语言。
\end{enumerate}

\paragraph{这正是开篇许下的承诺。}
这条“几何 $\to$ 概率 $\to$ 学习 $\to$ 智能”的弧线，并非事后追认。它在两处早已写明：其一，本书\emph{前言}立的旗——“沿‘经典模型 $\to$ 学习方法 $\to$ 具身智能’的脉络，架起一座由传统机器人学通向现代具身智能的桥梁”；其二，\cref{ch:intro} 借 Cadena 等\cite{cadena2016past} 勾出的 SLAM \emph{三个时代}（\cref{tab:intro-eras}）：古典时代奠定概率表述、算法分析时代吃透稀疏性与可观性、而我们身处的\emph{鲁棒感知时代}所追求的“超越基础几何的高层理解、任务驱动的感知”，恰恰就是这一部四章正在兑现的东西。绪论里那节“前沿与展望”（\cref{sec:intro-frontier}）当时只敢说“本书不会逐一展开”，把动态、可变形、协作等方向作为\emph{路标}先行预告；走到这里，路标后面的路，我们已经一节一节铺完了。

\paragraph{同一条弧线，还可以从“能力”这一侧读。}
上面四步是按\emph{方法}（几何、概率、学习、开放）切分的；换一个视角，从地图\emph{能表达到什么程度}切分，会得到一条几乎重合的阶梯。业界曾把视觉 SLAM 的能力粗分为三级\cite{davison2026computational}：\textbf{一级}只求稀疏定位（知道相机在哪、几个特征点在哪）；\textbf{二级}进一步建出\emph{稠密}几何（每个像素都有深度、表面被重建出来）；\textbf{三级}再叠上\emph{语义}（地图元素不只“在哪”，还“是什么”）；而这条阶梯之外、隐约可见的终点，是把场景组织成带物体实例、可被物理模拟的\emph{场景图}。这级别划分与本书的四步弧线并非巧合地吻合：一级对应经典几何与概率估计撑起的定位地基，二级是学习赋能稠密化的疆域（\cref{ch:learning,ch:dense}），三级正是度量-语义（\cref{ch:semantic}），而面向物体与开放词汇的场景图，恰是开放世界这一章（\cref{ch:openworld}）所指的方向。两种切法——按方法、按能力——殊途而同归，正说明这条弧线不是本书强加的叙事，而是这一领域\emph{自身}的演化纹理。

\paragraph{但要小心：层级不是“往上摞”。}
有一处易被误读：能力分级容易让人以为新一级是往一个已跑通的系统上\emph{加}一层。实情往往相反\cite{davison2026computational}——每一次表示的升级，都\emph{逼着整个系统重新设计}：一旦有了稠密地图，相机定位就该改用基于稠密模型的直接跟踪；一旦有了语义，稠密重建本身又能被语义反过来改善。实时系统的每个部件都在一个闭环里牵动其余部件，这与下一条洞见说的“层层包裹”是同一件事的两种说法——包裹，不是叠加。

\begin{insight}[弧线的形状：不是替换，而是层层包裹]
回望这四步，最该记住的不是“新方法淘汰了旧方法”，而是\emph{相反}：概率估计把几何放进后验里、学习把几何放进可微的一层里、开放世界把几何与语义放进同一张地图里。每一次“前沿”，都不是另起炉灶，而是给\emph{同一个内核}套上更大的外壳。这也是为什么本书坚持\emph{按主题而非按时代}组织——若把“经典”与“学习”割成两本不相干的书，读者就会错过这条弧线最要紧的信息：它们解的是\emph{同一个}问题。
\end{insight}

% ════════════════════════════════════════════════════════════════
\section{贯穿始终的几个思想}
\label{sec:epilogue-threads}

技术细节会过时——今天屠榜的系统，三五年后多半被新架构取代。但有几个思想，是无论前端换成什么网络、地图换成什么表示都\emph{不会}变的。它们是本书真正想留给读者的东西，也是这一章愿意反复重述的几条主线。

\paragraph{其一：一张因子图，一条母方程。}
本书从 \cref{ch:intro} 推出 MAP 母方程（\cref{eq:intro-nls}）
\begin{equation*}
  X^\star=\arg\min_X\ \sum_i \big\|\mathbf{e}_i(\mathbf{x}_i)\big\|_{\boldsymbol\Sigma_i^{-1}}^2
\end{equation*}
之后，它就再没真正离开过视野。视觉里程计的重投影误差是它的因子，回环是它的因子，IMU 预积分是它的因子（\cref{ch:vio}）；到了这一部，学习式光流是它的因子（\cref{ch:learning}）、动态物体的恒速运动是它的因子（\cref{ch:dynamic}）、物体与语义是它的因子（\cref{ch:semantic}）、开放词汇的特征观测仍是它的因子（\cref{ch:openworld}）。连 pointmap 这种“把流水线压进一次前向”的最激进做法，多图拼接时也要退回一个全局对齐的因子图。\textbf{“优化即推断”是这条主线的脊梁}：万变的是因子的来源（手工还是网络）、状态的表示（位姿点云还是点图），不变的是“把所有证据写成因子、在图上求一个一致的解”。读懂了这一点，再新的系统也只需问一句——\emph{它给这张图添了什么因子？}

\paragraph{其二：几何是免费而精确的先验，不要让网络重学它。}
深度学习最大的诱惑，是“端到端”——让一个网络从视频直接吐出轨迹。本书一以贯之地拒绝这条捷径，理由在 \cref{ch:learning} 反复出现：对极约束、刚体变换、投影几何，是\emph{精确且免费}的先验，硬要网络从数据里把它们“重新学一遍”，既浪费样本又难泛化。这不是保守，而是分工：把好学的（鲁棒对应、稠密深度、逐像素置信）交给网络，把该用约束的（多视一致、全局对齐）交给几何。差分式优化之所以重要，正因为它让这两者第一次可以\emph{端到端}地一起训练，而无须谁吞掉谁。还有一个更朴素、却几乎是决定性的论证支持这条分工\cite{davison2026computational}：设想真有一个网络，能一口吃下一百张图、直接吐出整个三维场景模型——那么当\emph{第一百零一张}图到来时该怎么办？把一百零一张图重新喂进网络、从头再算一遍，显然荒唐。只要承认一个系统要\emph{长期运行}、要把源源不断的新观测\emph{增量地}融进一个持久表示，就等于承认它需要概率状态估计那套“增量更新与融合”的工具——也就回到了本书的因子图与母方程。换言之，“增量融合”这一现实需求，本身就是几何式状态估计不会被端到端网络吞并的根由：它\emph{倒逼}出一个模块化、可持续更新的世界表示。

\paragraph{其三：对不确定性诚实。}
SLAM 与一般“感知”最深的区别，或许就在于它从不满足于一个点估计，而始终携带\emph{协方差}。一个回环到底可不可信、一块像素的颜色跨帧到底可不可靠、一个语义标签到底有多确定——这些“我有多确定”的问题，贯穿了鲁棒核（\cref{sec:nlopt-robust}）、抗外点的回环筛选（\cref{ch:loop}）、学习式光度不确定性（\cref{ch:learning}）乃至多类贝叶斯语义融合（\cref{ch:semantic}）。一个不诚实地报告自己不确定性的系统，在长时段运行里迟早会被一次过度自信的错误关联拖垮。把不确定性\emph{显式地建模、传播、并据以决策}，是 SLAM 留给整个机器人学的方法论遗产。

\paragraph{其四：清晰的记号，是发现错误的最好工具。}
这听上去不像一条“思想”，却是本书在每一章首都郑重登记记号、在前言里就申明“右扰动为主线、四元数取 Hamilton 约定、$\se(3)$ 平移在前”的良苦用心所在。当左右扰动混用、当帧的上下标方向不一致，错误往往不在算法而在\emph{记法}，且极难排查。把记号统一到底，让每一个 $\mathbf{T}_{ij}$、每一处求导都自洽，是一种廉价却极有回报的纪律——它让建模的错误在写下公式的那一刻就暴露出来，而不必等到系统跑飞。

\paragraph{其五：从相互矛盾的测量里萃取真相。}
一台机器人身上，相机、IMU、激光、轮速计各说各话，且都带噪声、偶尔撒谎。SLAM 的核心本领，正是把这些\emph{彼此冲突}的证据融成一个一致的、可被信任的世界估计——多传感器融合（\cref{ch:vio,ch:lidar_slam}）如此，抗感知混叠的鲁棒回环（\cref{sec:intro-need}）如此，把网络先验与几何约束调和到同一张图上（这一部）亦如此。这种“在冲突中求一致”的能力，是 SLAM 作为一门学问的真正难处，也是它的真正价值。

\begin{insight}[“SLAM 远未解决”——而这正是好消息]
时常有人断言“SLAM 已经解决了”。本书的立场与绪论一致（\cref{sec:intro-need}）：这个判断只对某个具体的\emph{机器人/环境/性能}三元组才有意义——二维室内、激光、厘米级精度，确可算基本解决；可一旦进入高速、强动态、长时段、开放语义、多机协同的设定，仍有大片基础问题敞开着。与其听信“已解决”，不如把目光投向当前技术水平之外两三步的问题——那里的工作，才经得起时间的考验。本章接下来要点的几个根本问题，正是这样的“两三步之外”。
\end{insight}

% ════════════════════════════════════════════════════════════════
\section{仍然敞开的根本问题}
\label{sec:epilogue-open}

把这一部读到底，最诚实的收获或许是：我们解决的每一个问题，都打开了更深的问题。下面几条不是习题，而是这条主线\emph{尚未走完}的方向——它们没有标准答案，留给读者，也留给这个领域。

\paragraph{机器人到底还需不需要“地图”？}
这是一个与绪论\emph{首尾相扣}的问题。\cref{sec:intro-need} 曾借 Cadena 等\cite{cadena2016past} 论证：只要还想要全局一致、还想抵御错误数据关联，回环——因而显式地图——就不可或缺。可到了基础模型时代，这个老问题被重新摆上了桌面（\cref{ch:openworld}）：当一个长上下文的视觉-语言模型可以直接“吃下”几十帧图像、凭其内部表示回答空间问题时，我们是否还需要一张\emph{外显}的、结构化的地图？目前的实验给出的是一个\emph{微妙}的答案而非定论——显式的几何结构与通用的学习策略，更像是\emph{互补}而非取代：结构带来可解释、可验证、可长期累积的一致性，策略带来开放与泛化。把二者如何结合得恰到好处，是这一部留下的最大悬念。

\paragraph{终身一致性：地图会“老”。}
本书绝大多数推导，默许了一个静态或缓变的世界。可真实世界是\emph{长期可变}的：家具被挪动、季节更替、建筑改建。一张昨天建好的地图，今天可能已处处失准。如何让一个系统\emph{终其一生}地维护一张始终一致的地图——既能容纳变化，又不被每一次变化推倒重来——是 \cref{ch:dynamic} 触及却远未穷尽的难题。它牵扯到一连串棘手的取舍：什么该记住、什么该遗忘、边缘化掉的历史一旦“烙死”便再难更改。

\paragraph{把不确定性贯彻到学习时代。}
经典后端给出的协方差有清晰的统计含义；可一个网络吐出的“置信图”究竟意味着什么？它是\emph{校准}的吗？当几何与学习被缝进同一张可微的图，如何让网络的不确定性与几何的协方差\emph{在同一尺度上对话}，使整个系统在分布外的场景里仍能诚实地说“我不确定”——这是把上一节“对不确定性诚实”这条主线推进到学习时代时，绕不开的开放问题。

\paragraph{鲁棒性、效率与保证。}
SLAM 之所以是机器人自主性的脊梁，正因它有\emph{严格的数学表述}。这份严格的潜力，远未被今天的系统兑现：我们仍缺少能在广义设定下，对解的\emph{次优性、计算效率与不确定性}同时给出\emph{保证}的方法。把“能跑的演示”提升为“有保证的系统”，把鲁棒性从“反复手工调参”中解放出来（\cref{tab:intro-eras} 所说“手工调参的诅咒”），仍是这条主线最艰难、也最值得的目标。

% ════════════════════════════════════════════════════════════════
\section{从空间感知到空间智能}
\label{sec:epilogue-toward}

\paragraph{先把“空间智能”说清楚。}
这一部反复出现“空间 AI”（Spatial AI）一词。\cref{ch:openworld} 已给过它一个\emph{偏语义}的工作内涵——地图超越纯几何、进而编码语义、支持高层推理。这里再补一个\emph{更工程化}的刻画，与之互补\cite{davison2026computational}：所谓空间 AI 系统，是一个\emph{实时}运行、主要靠视觉输入、\emph{增量地}建起并维护一份“近乎度量、至少局部如此、且人可理解”的场景表示、同时估计自身在其中位置、并能用\emph{少数性能指标}来量化其好坏的模块。这个刻画背后压着两条朴素却有力的假设。\textbf{其一}，对一台要长期运行、要做许多事先未必都已知的任务、还要与人交流的设备来说，它该维护的不是某个为单一任务定制的黑箱表示，而是一份\emph{通用、持久、贴近度量几何、且人能看懂}的世界模型——这恰是本书一路坚持“显式、结构化地图”的另一种说法。\textbf{其二}，五花八门的应用（送货无人机、家务机器人、AR 眼镜）所需的空间 AI 系统，其差别可以被\emph{相当少的几个性能参数}刻画清楚：除了人们熟知的全局定位精度与更新延迟，更贴近任务的还有“到表面接触的距离预测精度”“物体识别精度”“跟踪鲁棒性”这类指标。把要解的问题先用这样一个定义与一组指标钉住，本书后端那条“优化即推断”的主线，才有清晰的服务对象。

\paragraph{顺带澄清一个用词：“世界模型”。}
近年机器学习界热议的“世界模型”（world model），与本书一路在讲的\emph{场景表示}，所指其实是同一件事\cite{davison2026computational}：一份存在内存里、随新观测持续改写、也可主动遗忘以求高效的、关于设备与其环境状态的持久表示。差别只在出发点——许多端到端网络是\emph{无状态}的，吃进输入直接吐出输出，于是“给网络配一块持久记忆”被当成新课题；而 SLAM 这门学问，恰恰从第一天起就在做“持久表示”这件事。把这层窗户纸捅破有现实意义：它提醒我们，“世界模型”不必是一个抽象黑箱，本书所偏爱的——\emph{显式记下几何与物体}的表示，在效率、可组合、可解释、可长期累积上自有其不可替代的好处。

\paragraph{怎么衡量“做得好”：少信榜单，多信现场。}
还有一个贯穿这一部、却少被明说的态度问题：拿什么衡量一个 SLAM/空间 AI 系统的进步？计算机视觉惯于刷榜单，可单一数据集往往只测了\emph{精度}一个侧面，对同样要紧的\emph{鲁棒性、效率、灵活性}几乎无能为力，研究者还常只挑系统“跑得动”的那几段序列来报成绩\cite{davison2026computational}。SLAM 因其\emph{实时}本性、输出多样、对不同应用各有不同的“够用”标准，尤其难被一两个榜单刻画。事实上这一领域真正的里程碑，往往不是某个榜单上的零点几个百分点，而是一串可被人当场把玩的开源实时系统——本书一路点过的 MonoSLAM、PTAM、KinectFusion、LSD-SLAM、ORB-SLAM、DROID-SLAM、高斯泼溅 SLAM……都是这种“现场即证据”的标杆。这并非说基准毫无价值，而是提醒：把“能跑的演示”误当成“有保证的系统”固然危险（这正是上一节的开放问题），把“榜单分数”误当成“真实有用”同样危险。值得一提的是，面向真正具身的设备，未来最该被纳入的指标里，会出现一个今天罕见的维度——\textbf{片上数据搬运量}，以“比特 $\times$ 毫米”计：因为对一台要塞进眼镜、靠电池供电的设备而言，把数据从一处搬到另一处所耗的能量，往往比计算本身更昂贵。这条度量把我们引向一个纯算法视角看不见、却决定成败的层面——\emph{硬件}；本书在 \cref{sec:est-gbp} 已就此点到为止。

\paragraph{为什么“建好图”只是半个故事。}
至此，我们已能让机器人精确地知道\emph{自己在哪、世界长什么样、世界里有什么}。可这一切——定位、建图、语义——本身并不是目的。一台只会感知与建图、却不会\emph{行动}的机器人，与一台精密的测绘仪并无本质区别。机器人之所以是机器人，是因为它要\emph{进入}它所感知的世界、在其中\emph{做出决策}、并\emph{施加行动}去改变它。地图存在的全部理由，最终落在一句话上：\textbf{它是为了支撑决策与行动}。绪论坚持“SLAM 仍有价值”的第三条理由——许多任务显式或隐式地\emph{要求}一张全局一致的地图（\cref{sec:intro-need}）——指向的正是这件事：地图是任务的地基，而任务，意味着行动。

\paragraph{下一程：规划与控制。}
于是全书的重心，自然地从\emph{感知与空间智能}转向\emph{决策与行动}。有了地图与自身位姿，机器人接下来要回答两个层层递进的问题：

\begin{itemize}
  \item \emph{该往哪走？}——在已知（或正在建立）的地图上，找一条从此处到目标、既避开障碍又满足机器人运动学/动力学约束的路径乃至轨迹。这是\textbf{运动规划}的疆域（\cref{ch:planning}）：构型空间、图搜索、采样式方法、轨迹优化，皆为此而生。值得一提的是，规划与感知并非单向——一条好的路径会主动\emph{去往信息丰富处}以改善地图，这“任务驱动的感知”正是 \cref{ch:openworld} 与 \cref{tab:intro-eras} 鲁棒感知时代所点的方向，也让本书前后两半在此交汇。
  \item \emph{怎样走稳？}——把规划出的轨迹，变成驱动器上一刻不停的力与力矩，并在扰动、模型误差与噪声下\emph{稳定地}跟踪它。这是\textbf{控制}的疆域（\cref{ch:control}）：反馈与稳定性、最优控制、模型预测控制、以及面向机器人本体的动力学控制。
\end{itemize}

\noindent 这两步，与我们刚刚走完的感知主线共享着同一套数学底色——状态、不确定性、优化、流形上的运动——只是把它们从“估计世界”转向“在世界中行动”。\cref{ch:planning,ch:control} 将接续展开。至此，“感知—建图”这半张拼图已经拼好；接下来，是“决策—行动”那半张。机器人的自主，正是这两半合拢的地方。

\begin{insight}[全书的合拢：估计、规划、控制是同一门学问的三个侧面]
若说本书有一个贯穿\emph{始终}的隐线，那就是：无论估计（“世界是什么样”）、规划（“该怎么走”）还是控制（“怎样走稳”），底层都是\emph{在不确定性下、于流形上、求解一个带约束的优化问题}。SLAM 教我们把“理解世界”写成优化；规划教我们把“选择路径”写成优化；控制教我们把“稳定跟踪”写成优化。读者读到这里若有一种“似曾相识”的预感，那不是错觉——正是这种数学语言的\emph{统一}，让机器人学得以成为一门自洽的学问，而非一堆互不相干的技巧。带着这条隐线翻开下一部，会看得格外清楚。
\end{insight}

% ════════════════════════════════════════════════════════════════
\section*{几个值得带走的问题}
留几道开放的思考题，作为这一部、也作为整条 SLAM 主线的收束。它们都没有标准答案。

\begin{enumerate}
  \item \textbf{一张图，串起全书。}\;任取本书讲过的四个 SLAM 系统/方法（例如：经典 ORB-SLAM 的重投影 BA、VIO 的 IMU 预积分、DROID-SLAM 的光流因子、语义 SLAM 的物体因子），分别说出它给 \cref{eq:intro-nls} 这张因子图\emph{添了什么因子}、其残差与权重各来自手工还是网络。借此论证“同一条母方程贯穿始终”这一说法。
  \item \textbf{弧线的下一步。}\;本章把 SLAM 主线概括为“经典几何 $\to$ 概率估计 $\to$ 学习赋能 $\to$ 开放世界”。如果要你续上\emph{第五个}驿站，你认为它会是什么？给出一条理由，并说明它会给上面那张因子图带来什么新的因子或新的状态表示。
  \item \textbf{还需不需要地图？}\;就 \cref{sec:epilogue-open} 的第一个问题，选一个具体的机器人任务（如家庭服务、仓储分拣、灾后搜救），论证在这个任务里“外显结构化地图”与“长上下文模型的内部表示”应当如何分工。把你的论证回扣到 \cref{sec:intro-need} 的“机器人/环境/性能”三元组。
  \item \textbf{对不确定性诚实，到了学习时代。}\;经典后端的协方差有明确统计含义，网络吐出的“置信”却未必校准。设计一个\emph{思想实验}：如何检验一个学习式 SLAM 系统的不确定性是否“诚实”？（提示：考虑分布外输入下，它报告的置信与实际误差是否一致。）
  \item \textbf{从感知到行动的接口。}\;承接 \cref{sec:epilogue-toward}：一张带语义的地图，要供\emph{运动规划}（\cref{ch:planning}）使用，至少需要从中读出哪几样东西？再供\emph{控制}（\cref{ch:control}）使用呢？借此说清“地图是为决策与行动服务”这句话的具体含义。
\end{enumerate}

% ════════════════════════════════════════════════════════════════
\section*{延伸阅读}
\begin{itemize}
  \item \textbf{本章精神之源}：SLAM Handbook\cite{carlone2026handbook} 的结语——若干资深 SLAM 学者对这一领域的反思与寄语。本书未照搬其格言，而是择其与全书主线契合者（信任数学、清晰记号、SLAM 远未解决、从冲突测量中萃取真相、显式结构与学习策略互补），熔成本书自己的收束之辞。
  \item \textbf{空间 AI 的愿景蓝图}：Davison 在 SLAM Handbook 中关于\emph{空间 AI 计算结构}的专章\cite{davison2026computational}，把“SLAM 终将演化为空间 AI”这一判断展开为一整套前瞻——空间智能的能力阶梯、状态估计与机器学习如何长期共处、以及算法的图结构如何与未来硬件相匹配。本章对“几何不被端到端取代”“能力分层即本书弧线”等论断的综述，事实取自该章而立场熔成本书自己的口吻；其硬件与分布式算法一面，另见 \cref{sec:est-gbp} 对高斯置信传播的前瞻讨论。
  \item \textbf{首尾呼应}：本章多处回扣 \cref{ch:intro}——尤其 SLAM 三个时代（\cref{sec:intro-history,tab:intro-eras}）、“是否还需要 SLAM/地图”（\cref{sec:intro-need}）与“前沿与展望”（\cref{sec:intro-frontier}）。建议与绪论对读，体会这条弧线如何从“预告”走到“兑现”。其历史叙事的源头仍是 Cadena 等\cite{cadena2016past}。
  \item \textbf{承上}：本部四章——学习赋能（\cref{ch:learning}）、动态与可变形（\cref{ch:dynamic}）、度量-语义（\cref{ch:semantic}）、开放世界空间 AI（\cref{ch:openworld}）——是本章回望的对象。神经渲染（NeRF/3DGS 一系）作为稠密表示的一条新分支，在 \cref{ch:learning,ch:dense} 已指明其归属。
  \item \textbf{启下}：运动规划（\cref{ch:planning}）与控制（\cref{ch:control}）承接“从感知到行动”的转折，是全书后半程的起点。
\end{itemize}

\section*{本章与全书的关系}
\begin{table}[H]\centering\small
\begin{tabular}{lll}
\toprule
章节 & 关系 & 衔接点 \\
\midrule
\cref{ch:intro} 初识 SLAM & 首尾呼应 & 三时代弧线、母方程、“是否需要地图”在此呼应闭环 \\
\cref{ch:learning} 学习赋能 & 回望（承上）& 几何被学习包裹、不被取代 \\
\cref{ch:dynamic} 动态可变形 & 回望（承上）& “地图会老”的终身一致性开放问题 \\
\cref{ch:semantic} 度量-语义 & 回望（承上）& 地图从“记几何”到“记是什么” \\
\cref{ch:openworld} 开放世界 & 回望（承上）& 开放词汇、“是否还需要外显地图” \\
\cref{ch:planning} 运动规划 & 过渡（启下）& 地图为决策服务：该往哪走 \\
\cref{ch:control} 控制 & 过渡（启下）& 把轨迹变成稳定的行动：怎样走稳 \\
\bottomrule
\end{tabular}
\end{table}
